\section{What Actually Exists}

This is the first of four fronts on which the objection has to be met, and it is the one that carries the others. If particulars are metaphysically second-rate---if to be this one rather than that one is to be a diminished instance of something better---then every later section of this article is built on sand, and the objectors were right for reasons that have nothing to do with theology. So the case against the particular has to be stated at its strongest before anything is said in its favour, and it has to be stated in the philosophers who made it, not in the popular version.

\subsection{The case against the particular}

Two claims stand behind the objection, and they come from the two men who founded the discipline.

The first is Plato's, and it is the more famous. Sensible particulars, on the account the \emph{Phaedo} gives, do not merely differ from the Forms; they fall short of them, and they are known to fall short. Socrates asks whether equal sticks and stones are equal ``in the same sense in which absolute equality is equal? or do they fall short of this perfect equality in a measure?''---and gets the answer the argument needs: ``Yes, he said, in a very great measure too.'' The observation is then generalized: when anyone looking at an object ``observes that the thing which he sees aims at being some other thing, but falls short of, and cannot be, that other thing, but is inferior,'' he must already have known the thing it fails to be.\endnote{Plato, \emph{Phaedo} 74a--75b, in Benjamin Jowett's public-domain English translation (3rd~ed., Oxford, 1892); the whole dialogue downloaded and read 26 July 2026 in the Project Gutenberg text, ebook 1658. Jowett died in 1893 and the translation is out of copyright by age. The Stephanus numbers are supplied from the standard divisions of the dialogue and are not printed in the witness read; the passage is identified by its wording. Nothing here adjudicates the vexed question of how far the historical Socrates held the theory the dialogue gives him.} The particular is thus defined by a deficit. It is an approximation to what it would be if it could.

The \emph{Republic} draws the epistemological consequence, and it is sharper than the metaphysics. Those who love sights and sounds---who are ``fond of fine tones and colours and forms and all the artificial products that are made out of them,'' but whose ``mind is incapable of seeing or loving absolute beauty''---are not merely less cultivated. They are asleep. Knowledge, Plato argues there, ``is relative to being and knows being''; what stands between being and not-being is the object not of knowledge but of opinion, and the many beautiful things are exactly what stands between.\endnote{Plato, \emph{Republic} V, 475d--480a, Jowett translation; downloaded and read 26 July 2026 in the Project Gutenberg text, ebook 1497. The dream-and-waking contrast immediately follows the sentence quoted. The article uses the passage to state the objection's classical credentials and does not undertake an interpretation of Plato's theory of Forms, on which the literature is enormous and the article takes no position.} That is the objection's ancestral form, and it is worth seeing how much stronger it is than the theological version. It does not say that God would find particulars beneath him. It says that particulars are, strictly, not fully knowable---and if that is right, then a revelation given in particulars is a revelation given in the medium least capable of carrying truth.

The second claim is Aristotle's, and it is harder, because Aristotle is on the other side of the Platonic question and reaches a version of the same difficulty anyway. He puts it as an aporia, in the form of a dilemma about the first principles of things. If the principles are universal, ``they will not be substances; for everything that is common indicates not a \emph{this} but a \emph{such}, but substance is a \emph{this}.'' If instead they are individual, ``they will not be knowable; for the knowledge of anything is universal.''\endnote{Aristotle, \emph{Metaphysics} III (B), in W.~D. Ross's English translation (first published Oxford, 1908); the complete translation downloaded and read 26 July 2026 in the transcription at the Internet Classics Archive. Ross's translation is out of copyright by age. The Oxford printing itself was not collated: a scan of the 1928 second edition was obtained at the Internet Archive, but its optical transcription proved too degraded to check these sentences against, and that limit is recorded in the terminal appendix and the research record. The witness read prints book and chapter divisions but no Bekker numbers; the sentences are identified by wording.} The dilemma is restated in Book XIII with the second horn made explicit: individual principles ``are not knowable; for they are not universal, and knowledge is of universals.''

Then Aristotle presses his own side of it further than any opponent of particularity has since. In Book VII he argues that ``no universal attribute is a substance, and this is plain also from the fact that no common predicate indicates a \emph{this}, but rather a \emph{such}''---so the universal cannot be what primarily exists. And in the same book he concedes the price: ``there is neither definition of nor demonstration about sensible individual substances, because they have matter whose nature is such that they are capable both of being and of not being.''\endnote{Aristotle, \emph{Metaphysics} VII (Z), chs.~13 and 15, Ross translation, same witness, read 26 July 2026. Chapter~8 supplies the formula the later argument uses: ``the whole \emph{this}, Callias or Socrates, is analogous to \emph{this brazen sphere}, but man and animal to \emph{brazen sphere} in general.''}

Put the two horns together and the objection reaches its full metaphysical strength, which is considerably beyond anything Celsus or Kant needed. It is not merely that the particular is humble. It is that what exists is not definable and what is definable does not exist---so that being and intelligibility come apart at exactly the point where the Christian claim wants them joined. A God who acts in particulars acts in what cannot be said; a God who is to be known must be known in universals, which are not the sort of thing that acts. That is the deepest form of the case this article has to answer, and no theological consideration whatever will touch it.

\subsection{Aristotle answers it himself}

The answer begins where the objection ends, and it is Aristotle's.

Having stated the difficulty about knowledge and universals as ``of all the points we have mentioned, the greatest difficulty,'' he does not leave it standing. He distinguishes:

\begin{quote}\small
yet the statement is in a sense true, although in a sense it is not. For knowledge, like the verb \emph{to know}, means two things, of which one is potential and one actual. The potency, being, as matter, universal and indefinite, deals with the universal and indefinite; but the actuality, being definite, deals with a definite object, being a \emph{this}, it deals with a \emph{this}.\endnote{Aristotle, \emph{Metaphysics} XIII (M), ch.~10, Ross translation, same witness, read 26 July 2026. The sentence immediately preceding states the difficulty in the form the objection uses: ``The statement that all knowledge is universal, so that the principles of things must also be universal and not separate substances, presents indeed, of all the points we have mentioned, the greatest difficulty \ldots'' The article's use of this passage---as the decisive concession that actual knowing terminates in a singular---is an interpretation, and a contested one in the literature on Book M, which this article does not survey.}
\end{quote}

That is the hinge of the entire metaphysical front, and it should be read slowly, because it comes from the objection's own founder. The premise the objector needs is that knowledge is of universals \emph{tout court}. What Aristotle grants is narrower: knowledge \emph{in potency} is of the universal, because a capacity, like matter, is indeterminate until it is exercised. Knowledge in act is of a determinate object---and a determinate object is a \emph{this}. Universality, on this account, is not the perfection of knowing. It is the condition of knowing before it has happened.

Everything the later argument needs is contained in that distinction. The objection assumed a scale running upward from the concrete to the abstract, with God at the top as the most general of all. Aristotle's analysis puts the abstract not above the concrete but before it, on the side of potency; and what stands at the top of that scale is not the most general but the most actual. The tradition that followed him drew exactly this conclusion about the divine intellect: God's knowledge does not run through the universal, because God's knowing is not a potency being exercised. The sequence is not an apologetic invention. It is what the premises give.

Two things follow that are worth marking before the argument goes on. The first is that Plato's claim survives in a weakened form and should be granted rather than fought: sensible particulars really are known by us imperfectly, and our concepts really are more stable than the things they are about. What does not follow is Plato's inference---that the deficiency lies in the particular rather than in the knower. The \emph{Phaedo}'s argument establishes that the equal sticks fall short of equality; it does not establish that they fall short of \emph{being}, and the two are quietly run together. The second is that Aristotle's concession in Book VII stands: individual sensible substances admit of no definition. That is a real cost, and this article does not evade it. Its answer is that indefinability is not unintelligibility---the concrete singular is what a definition is about and always more than the definition manages to say---and that the concession is in any case fatal to the objector's position rather than to this one, since a God who dealt only in definables would be dealing only in what does not exist.

\subsection{No one has ever met a universal}

Begin with an observation that requires no theology: everything that exists, exists as a particular. There is no such thing as \emph{dog} trotting across a road; there are only individual dogs. There is no \emph{humanity} that suffers or rejoices; there are only men and women, each with a fingerprint. Universals---dog, humanity, triangle, justice---are real objects of knowledge, and nothing here depends on denying their importance. But they do not walk about. In the classical account that this article adopts, a common nature has its concrete being only in individuals; taken apart from all individuals it exists, as Aquinas puts it, only ``as existing in an intellect either human or Divine.''\endnote{Thomas Aquinas, \emph{Summa theologiae} III, q.~4, a.~4, corpus, in the public-domain English Dominican translation, checked at the New Advent transcription, 24 July 2026. Aquinas is there rejecting both the Platonist subsistent universal and the suggestion that the Word could have assumed human nature ``abstracted from all individuals.'' The moderate-realist account of universals adopted in this section is a school position argued by the article, not a dogma; the article's later theological claims do not stand or fall with its technical formulation.}

This gives the objection its first jolt. The objector says: the universal is nobler than the particular, so a God who deals in particulars is slumming. But nobler \emph{as what}? As an object of contemplation, perhaps. As an agent, the universal is not merely humbler than the particular---it is nothing at all. Abstractions do not act. ``Humanity'' has never fed anyone; particular people have. Aquinas states the principle in passing while defining what a person is: rational substances ``are not only made to act, like others; but \ldots\ can act of themselves; for actions belong to singulars.''\endnote{Aquinas, \emph{Summa theologiae} I, q.~29, a.~1, corpus, English Dominican translation, checked at the New Advent transcription, 25 July 2026. The scholastic tag for the clause quoted is \emph{actiones sunt singularium}; it is mentioned here as the received label for the doctrine, not as a separately collated quotation.} If God is to \emph{do} anything in the world---create it, sustain it, rescue it---the doing will terminate in particulars, because there is nowhere else for action to land. A God who dealt only in universals would not be a more exalted God. He would be an idle one.

\subsection{Abstraction is our limitation, not God's grandeur}

The second jolt is sharper. Ask where abstraction comes from. In the Thomistic analysis of knowledge, it is the signature of a \emph{finite} intellect. Human minds cannot grasp the singular in its full concreteness directly; we know material singulars, Aquinas argues, only indirectly, because ``our intellect cannot know the singular in material things directly and primarily,'' knowing instead what it abstracts---and ``what is abstracted from individual matter is the universal.''\endnote{Aquinas, \emph{Summa theologiae} I, q.~86, a.~1, corpus, English Dominican translation, checked 24 July 2026. The claim concerns direct intellectual knowledge of \emph{material} singulars and includes Aquinas's account of indirect knowledge by reflection on phantasms; it is not a denial that human beings know individuals at all.} We trade the singular for the general because the singular, in its inexhaustible detail, exceeds us. Abstraction is a coping mechanism. It is a glorious one---all science lives in it---but it is the strategy of a mind too small to hold the individual whole.

God's knowledge runs in exactly the opposite direction. Because God is the cause of things, ``His knowledge extends as far as His causality extends''; and since his causality reaches not only forms but the very matter by which things are individuated, his knowledge reaches ``to singular things, which are individualized by matter.''\endnote{Aquinas, \emph{Summa theologiae} I, q.~14, a.~11, corpus, English Dominican translation, checked 24 July 2026.} Likewise providence: the causality of the first agent ``extends to all being, not only as to constituent principles of species, but also as to the individualizing principles,'' so that all things fall under providence ``not only in general, but even in their own individual selves.''\endnote{Aquinas, \emph{Summa theologiae} I, q.~22, a.~2, corpus, English Dominican translation, checked 24 July 2026. Aquinas's account of providence includes real secondary causes and divine permission; nothing here licenses reading particular sufferings as particular punishments.}

Here the objection quietly inverts. The philosopher who insists that God must relate to the world only through generalities---species, laws, averages---has not purified the idea of God. He has projected onto God the one intellectual limitation that is most distinctively ours. A God of universals only would be a God who knows the way a creature knows: by losing the individual in the type. The classical doctrine is stranger and better. Every sparrow is known singly, not as an instance of sparrowhood but as \emph{this} sparrow, because the act by which it exists at all is an act of God bearing on it in particular.

Origen had already made the point against Celsus, and in almost these words. Where Celsus had held that God ``takes care of the whole,'' Origen answers that God takes care ``not, as Celsus supposes, merely of the whole, but beyond the whole, in a special degree of every rational being.''\endnote{Origen, \emph{Contra Celsum} IV.99, Crombie translation, checked at the New Advent transcription, 25 July 2026. Origen's argument in the surrounding chapters that the world was made for the sake of rational creatures is his own, engages Stoic and Platonist cosmology, and is not adopted here; what is taken from the passage is only the distinction between care for the whole and care for each, which is the point at issue.} The contrast is exact, and it is the hinge of the whole dispute. Celsus's providence is distributive: it reaches individuals through the system. Origen's, and afterward Aquinas's, is direct: it reaches the system through individuals, because individuals are what there are.

\subsection{What makes this one this one}

If particulars are what exist, it is worth asking what makes a particular the particular it is---and the classical answer bears directly on why the Incarnation had to happen to somebody.

The question is a real one, because natures do not individuate themselves. Human nature is what Socrates and Plato have in common; it cannot also be what distinguishes them, or it would be doing contradictory work at the same moment. Aquinas's answer, taken over from Aristotle and sharpened, is that in material things the principle of individuation is matter---not matter in general, which is as common as the nature, but this matter, marked out here and now. The formulation is given most cleanly in his early treatise \emph{De ente et essentia}: matter is not the principle of individuation ``taken in just any way, but only designated matter. And I call designated matter that which is considered under determinate dimensions.''\endnote{Thomas Aquinas, \emph{De ente et essentia}, in the chapter treating essence in composite substances---ch.~2 in the common division, printed as ch.~1 of the treatise proper at the Corpus Thomisticum, which sets the prooemium apart and renumbers accordingly. Latin checked at the Corpus Thomisticum text, 25 July 2026: ``materia non quolibet modo accepta est individuationis principium, sed solum materia signata. Et dico materiam signatam, quae sub determinatis dimensionibus consideratur \ldots\ In diffinitione autem hominis ponitur materia non signata; non enim in diffinitione hominis ponitur hoc os et haec caro, sed os et caro absolute.'' The English here and in the next sentence is this article's own labeled working translation; no public-domain English translation of the treatise was identified for this study. The chapter-numbering divergence is recorded so that either edition can be checked.} And then the sentence that matters most for everything this article goes on to say: in the definition of man, non-designated matter is what is posited, ``for there is not posited in the definition of man \emph{this} bone and \emph{this} flesh, but bone and flesh absolutely.''

Read that slowly. The definition of man---the universal, the thing an intellect can hold---contains bone and flesh in general. It does not contain this bone. Nothing definable does. The concrete individual is precisely what falls outside every definition, which is why, as an objection in the \emph{Summa} puts it and Aquinas partly concedes, ``nothing singular can be subject to definition.''\endnote{Aquinas, \emph{Summa theologiae} I, q.~29, a.~1, objection~1 with its reply, English Dominican translation, checked 25 July 2026. The reply concedes that ``this or that singular may not be definable'' while holding that what belongs to the general idea of singularity can be defined. The concession is the point relied on here; the full reply is not being reduced to it.} The singular is not a deficient universal. It is what the universal is about, and it is always more than the universal manages to say.

The doctrine has a striking test case, and Aquinas applies it without flinching. If matter individuates, then beings without matter cannot be individuated within a species---so each angel must be its own species. He says exactly that: ``such things as agree in species but differ in number, agree in form, but are distinguished materially. If, therefore, the angels be not composed of matter and form \ldots\ it follows that it is impossible for two angels to be of one species; just as it would be impossible for there to be several whitenesses apart, or several humanities.''\endnote{Aquinas, \emph{Summa theologiae} I, q.~50, a.~4, corpus, English Dominican translation, checked 25 July 2026. The angelology is a school position contested in the schools---Scotus, among others, rejected it---and nothing in this article's theological claims depends on it. It is cited because it shows the individuation doctrine being applied consistently where the application costs something, which is a mark of a principle held rather than assumed.} Whatever one makes of the angelology, the method is instructive: the principle is pressed where it hurts, rather than quietly suspended.

It should be said that this was never the only account on offer. Duns Scotus, a generation later, held that the common nature is real and ``indifferent'' to existing in any number of individuals, but that in each individual it is contracted by a positive individuating principle proper to that individual---the \emph{haecceitas}, the thisness---so that the humanity of Socrates is made individual by Socrates's haecceity and the humanity of Plato by Plato's.\endnote{Reported at a stated secondary-witness ceiling from Thomas Williams, ``John Duns Scotus,'' \emph{Stanford Encyclopedia of Philosophy} (first published 31 May 2001; substantive revision 15 December 2025), checked 25 July 2026, which supplies the account summarized here and identifies the primary locus as \emph{Ordinatio} II, d.~3, pars~1, qq.~1--6. Scotus's text was not collated for this article, no Scotist wording is quoted anywhere in it, and no claim finer than the summary given---common nature, contraction, haecceity as a positive individuating principle---is made. The Thomist--Scotist dispute over individuation is not adjudicated here.} The disagreement between the schools is real, is not settled in these pages, and does not need to be, because the two accounts converge on exactly what this article requires. On either view the individual is not a smudged copy of the species. On either view individuality is a positive perfection rather than a defect in the universal---for Scotus explicitly, since haecceity is something added rather than something lacking; for Aquinas by consequence, since designated matter is what a nature actually has when it actually is. The objector assumes that to be this one rather than that one is to be diminished by a limit. Both schools deny it.

\subsection{The dignity of being someone}

The classical tradition goes further, and the further step is the one the objection never anticipates. It does not merely defend the particular; it locates the summit of created being there.

Aquinas builds the case in stages. Individuality belongs in a special way to substances, because a substance is individuated by itself while accidents are individuated by their subject---this whiteness is called ``this'' because it is in this thing. Then, he says, ``in a more special and perfect way, the particular and the individual are found in the rational substances which have dominion over their own actions.'' And so such individuals receive a name of their own: person.\endnote{Aquinas, \emph{Summa theologiae} I, q.~29, a.~1, corpus, English Dominican translation, checked 25 July 2026. The definition under discussion is Boethius's---``an individual substance of a rational nature''---and Aquinas's argument is an exposition of why each of its terms is required.} The conclusion, when it comes, is stated without hedging: ``\emph{Person} signifies what is most perfect in all nature---that is, a subsistent individual of a rational nature.''\endnote{Aquinas, \emph{Summa theologiae} I, q.~29, a.~3, corpus, English Dominican translation, checked 25 July 2026. The article asks whether ``person'' may be said of God, and answers that it may, since ``everything that is perfect must be attributed to God,'' though ``not, however, as it is applied to creatures, but in a more excellent way.'' That analogical qualification is essential and is preserved throughout this article: nothing here predicates individuality of God univocally with creatures, and the article's own terminology section restates the limit.}

That sentence should be set against the objection's whole aesthetic. The objector's scale of value runs upward from the concrete toward the abstract: this man, then humanity, then rational nature as such, then pure form, then God at the summit as the most general of all. The classical scale runs the other way at the decisive point. What is most perfect in all nature is not a nature. It is a subsistent individual---someone. Universality is not the direction of ascent; subsistence is.

The reason lies deeper still, in what the tradition says about existence itself. To be actual is not merely to instantiate a form; it is to have the act by which anything at all is. And that act, on the classical account, is what God gives---immediately, and continuously---to each thing that is. Aquinas draws the consequence in a passage that quietly dismantles the picture of a God who governs from a distance through general laws. God is in all things, he argues, not as part of their essence and not as an accident, but ``as an agent is present to that upon which it works''; since ``created being must be His proper effect,'' and since God causes that effect ``not only when they first begin to be, but as long as they are preserved in being,'' God must be present to each thing for as long as it exists. Then the clinching step: ``being is innermost in each thing and most fundamentally inherent in all things \ldots\ Hence it must be that God is in all things, and innermostly.''\endnote{Aquinas, \emph{Summa theologiae} I, q.~8, a.~1, corpus, English Dominican translation, checked 25 July 2026. The third reply of the same article---``it belongs to the great power of God that He acts immediately in all things. Hence nothing is distant from Him''---bears directly on the objection's picture of a remote God dealing in generalities. The presence asserted is that of causal immediacy, not of composition; Aquinas denies composition expressly at I, q.~3, a.~8, cited above.}

The metaphysical distance between God and a general law could hardly be greater. A law abstracts from the individual: that is what makes it a law, and it is why a law can never say anyone's name. The act of being does the opposite---it constitutes the individual as the individual it is. The objector's proposal, that God relate to creatures as a statute relates to citizens, therefore does not exalt God above particular involvement. It substitutes, for the most intimate relation there could be, one of the most external available.

Two doctrines standing behind that paragraph are argued elsewhere rather than here, and it is worth saying so plainly rather than smuggling them in. That no finite nature accounts for its own existing, so that a creature's being is received and participated rather than owned, is the burden of a companion study in this collection; so is the account of how perfections found in creatures may be predicated of their source analogically, neither univocally nor equivocally.\endnote{The essence--existence distinction, participated \emph{esse}, and the analogical predication of creaturely perfections of God are argued in the companion studies of this collection---particularly the first part of \emph{The Creature Before God}---and are not re-argued here. This article borrows only the consequence stated in the text: that what a creature most fundamentally receives is its act of being, and that this act is received by an individual. Readers who reject the participation framework can still take the present section's argument in its weaker form, which needs only that abstraction is a finite knower's device and that action terminates in singulars.} Neither is re-argued in these pages. What this section borrows from them is a single consequence: if creatures participate in existence rather than possessing it natively, then what a creature most fundamentally receives is not its species but its act of being---and that act is received by \emph{this} thing. Participation is accordingly not a general relation obtaining between two categories. It is a relation each singular has, one at a time, to the one who is.

\subsection{The universal cause is not a general thing}

One confusion remains to be cleared, because the objection lives off it. ``God is universal'' is true in one sense and false in another. God is the universal \emph{cause}: nothing exists that does not receive its being from him. But a universal cause is not a general \emph{thing}, a widest category, a cosmic gas. Aquinas is blunt: God is not in any genus at all.\endnote{Aquinas, \emph{Summa theologiae} I, q.~3, a.~5 (God is not in a genus) and a.~8 (God does not enter into composition with things), English Dominican translation, checked 24 July 2026.} He is not the most abstract item on the shelf; he is not on the shelf. The choice the objector offers---either an absolute God aloof from particulars, or a tribal god entangled in them---is therefore a false alternative resting on a class error. Precisely because God is not one being among beings, his relation to the particular is not competition for space within it. The universal cause is nearer to each singular thing than any general law could be, for the law abstracts from the individual and the cause gives the individual its very act of existing.

The same distinction disposes of Celsus's dilemma about descent. He argued that if God comes down among men he must change, since to enter the particular is to be altered by it. That follows only if God is the sort of thing that occupies a position, so that arriving somewhere entails leaving somewhere else. Origen saw the point and made it first: the Word ``does not give His place or vacate His own seat, so that one place should be empty of Him, and another which did not formerly contain Him be filled''---for ``in speaking of His quitting one place and occupying another, we do not mean such expressions to be taken topically.''\endnote{Origen, \emph{Contra Celsum} IV.5, Crombie translation, checked at the New Advent transcription, 25 July 2026. Origen appeals to Jeremiah's ``Do I not fill heaven and earth?'' and to Paul's Areopagus speech. What the passage settles is only the coherence question---that presence in a particular need not be desertion of the whole. The properly Christological questions raised by the Incarnation are not settled by it and are treated at their own place below.}

A third position should be recorded, because it changes what the dispute is about rather than taking a side in it. The nominalist tradition that runs through William of Ockham denies that there are universal entities outside the mind at all: on the reported account, the only real universals are universal concepts, and derivatively the spoken or written terms expressing them, so that Ockham's metaphysics is an ontological reduction which ``emphatically'' denies universal, quantitative, and relational entities.\endnote{Reported at a stated secondary-witness ceiling from the \emph{Stanford Encyclopedia of Philosophy} entry ``William of Ockham'' (first published 16 August 2002; substantive revision 11 September 2024), downloaded and read 26 July 2026. No text of Ockham was collated for this article and no Ockhamist wording is quoted anywhere in it. The same entry records a point worth passing on, since it is often got wrong: the slogan usually called Ockham's Razor, that entities must not be multiplied beyond necessity, is nowhere to be found in Ockham's texts, though the sentiment is his.} If that is right, then the question the Thomist and the Scotist were arguing about---what contracts a common nature to this individual---has no object, because there is no common nature in things to be contracted. The consequence for the present argument is worth stating plainly. On nominalist premises the article's first claim is not weakened but made trivial: everything that exists is a particular, because there is nothing else for anything to be. The three schools thus disagree about the analysis and agree about the datum, and the datum is all this article uses. That is why the section can afford to leave the dispute open. A conclusion that follows on moderate realism, on Scotism, and on nominalism alike does not need the dispute settled; it needs only that no one of the three supplies the objector with the premise he requires, and none of them does.

\begin{studybox}{Project synthesis --- the metaphysical front}
\textbf{Judgment.} The front is won, and it is won on the objector's own ground. Particulars are not metaphysically second-rate: what exists concretely is singular, action terminates in singulars, individuality is a positive perfection rather than a limit, and the summit of created being is a subsistent individual and not a nature. One genuine concession is owed and is made below rather than evaded. The theological sections of this article therefore rest on something, and not on an apologetic preference for the concrete.

\textbf{Reasons.} (1)~The strongest form of the case against the particular is not theological but Aristotle's aporia---that what exists is not definable and what is definable does not exist---and Aristotle himself resolves it, by distinguishing knowledge in potency, which is of the universal and indefinite, from knowledge in act, which ``being a \emph{this}, deals with a \emph{this}.'' Universality is thus the condition of knowing before it has happened, not its perfection. (2)~Plato's argument, read at its own locus, establishes that sensible things fall short of the Forms as measures; it does not establish that they fall short of \emph{being}, and the objection needs the second. (3)~Abstraction is on the classical analysis the signature of a finite intellect, so a God confined to generalities would be modelled on a limitation of ours. (4)~The three live accounts of individuation---Thomist designated matter, Scotist haecceity, Ockhamist denial that the question arises---converge on the one premise the argument uses, that the individual is not a deficient copy of a species. An objector must therefore defeat all three, and no historical objector has attempted it.

\textbf{Strongest counterargument.} That the concession is larger than the section admits. Aristotle grants, and this article grants with him, that there is ``neither definition of nor demonstration about sensible individual substances.'' If the singular is strictly indemonstrable, then a revelation delivered in singulars is delivered in a medium about which nothing can be proved---and that is not a quibble, since it is precisely what the historical front will press. Worse, the resolution leaned on is a single passage of \emph{Metaphysics} M, whose interpretation is contested in the literature this article has not surveyed; a reader entitled to say that Ross's Aristotle is being read tendentiously is entitled to say that the hinge is unsupported. The honest position is that the metaphysical front is won against the claim that particulars are inferior beings, and is \emph{not} won against the claim that particulars are epistemically intractable. That second claim is conceded here and answered, if at all, in the section on the ditch.

\textbf{What would change this judgment.} (1)~A demonstration that the potency/act distinction in \emph{Metaphysics} M cannot bear the reading given, either from the Greek or from the run of Aristotle's argument in Books M--N. The section's hinge would go, and with it the claim that the objection's own founder concedes the point. (2)~A defensible account of individuation on which individuality \emph{is} a privation---something like a rigorously worked out Neoplatonism in which the singular is constituted by the failure of a form to be wholly received. All three live accounts deny this; a fourth that established it would restore the objection at full strength. (3)~A demonstration that ``actions belong to singulars'' is false, that is, that there are agents which are not individuals. That would remove the argument that a God who acts must act among particulars, which is the load-bearing member of everything that follows.
\end{studybox}

None of this yet proves that God chose Abraham or took flesh from Mary. Metaphysics cannot deliver those claims, and this article will not pretend that it can. What the metaphysics does is demolish the veto. There is no rational principle by which the concrete is unworthy of God. If anything, the analysis points the other way: a God who knows and sustains every singular directly, and for whom being someone is the highest thing there is in creation, is a God for whom addressing \emph{someone in particular} is not an embarrassment but the obvious idiom. Whether he has actually spoken is a question of fact, not of fittingness---and questions of fact are settled, if at all, in history. To history, then.
